% !TEX root = ../Main.tex
\chapter{运动状态的改变}

运动状态由\textbf{速度} $v$ 刻画，而速度是矢量：
\[
v\ \begin{cases}\text{大小}\\[2pt]\text{方向}\end{cases}
\]
所以"运动状态改变"指的是速度的大小\textbf{或}方向发生了变化——接着上一章那句话，二者只要变一个，状态就变了。

\section{两个例子}

\ex{匀速圆周运动。}

\sol{
    速度大小始终不变，但方向时刻在变（永远沿切线）。方向变了，速度这个矢量就变了，所以匀速圆周运动\textbf{运动状态一直在改变}。

    \begin{center}
    \begin{tikzpicture}[>=Stealth, scale=1]
        \draw (0,0) circle (1);
        \filldraw (1,0) circle (1.2pt);
        \draw[->, thick, red] (1,0) -- (1,0.9) node[right] {$v$};
        \filldraw (0,1) circle (1.2pt);
        \draw[->, thick, red] (0,1) -- (-0.9,1) node[above] {$v$};
        \filldraw (-1,0) circle (1.2pt);
        \draw[->, thick, red] (-1,0) -- (-1,-0.9) node[left] {$v$};
    \end{tikzpicture}
    \end{center}
}

\ex{竖直上抛运动的最高点。}

\sol{
    到最高点时速度恰好为零，但此刻物体仍受重力，下一瞬就开始往下加速。所以"速度为零"\textbf{不等于}"运动状态不变"——最高点处运动状态照样在改变。

    \begin{center}
    \begin{tikzpicture}[>=Stealth, scale=1]
        \draw[->] (0,0) -- (0,2.6) node[above] {};
        \draw[thick, blue] (-0.5,0) .. controls (-0.15,2.2) and (0.15,2.2) .. (0.5,0);
        \filldraw (0,2.15) circle (1.2pt) node[above] {\footnotesize 最高点};
        \node at (-0.9,1.1) {$v$};
    \end{tikzpicture}
    \end{center}
}

\section{判定：靠力还是靠速度}

\state{力是改变运动状态的原因}{
    \textbf{力是改变物体运动状态的原因}。由这条原理，判断"运动状态变没变"有两套等价说法：
    \begin{enumerate}
        \item 看受力：不受力，或所受合力为 $0$ $\Rightarrow$ 运动状态不变；
        \item 看速度：速度 $v$（大小和方向）不变 $\Rightarrow$ 运动状态不变。
    \end{enumerate}
    两个角度结论一致：合力为零 $\Leftrightarrow$ $v$ 不变。
}
